% PAGES: 147
% Continues sec04_c3_3.tex (pages 145-146), which left the \begin{proof}
% of Lemma 4.11 open right after equation (4.52). This file completes
% that same proof (no re-opening) -- it concludes with the printed square
% near the top of page 147 (folio 131), which also closes section 4.5.
%
% Section 4.6 "References" then opens and runs to the bottom of page 147
% (four reference paragraphs, the last one ending "...Lax and Zalcman
% \cite{27}."). No environment is left open at the end of this file, and
% no sentence, \[ \], or begin{...} is left unclosed. This is the last
% page of the c3 chunk (PDF pages 141-147 / folios 125-131); the next
% chunk (c4, pages 148-154) should open fresh, not as a continuation of
% anything in this file.

Recall from (4.43) that the matrix $B_N$ has the following $N$ eigenvalues:
\[
\mu_j \;=\; 2 - 2\cos\left(\frac{j\pi}{N+1}\right), \quad \text{for}\ j = 1:N.
\]
Then, from (4.52), it follows that
\begin{equation}
\begin{split}
\lambda_j &\;=\; d - 2a + a\mu_j \;=\; d - a(2-\mu_j)\\
&\;=\; d - 2a\cos\left(\frac{\pi j}{N+1}\right).
\end{split}
\end{equation}
We conclude that, for $j = 1:N$, $\lambda_j$ given by (4.53) are eigenvalues of the
matrix $T_N$ with corresponding eigenvectors $v_j$ given by (4.44).

The $N \times N$ matrix has exactly $N$ eigenvalues, counted with their
multiplicities; cf. Theorem 4.2. We conclude that $\lambda_j$, $j = 1:N$, are all the
eigenvalues of $T_N$.
\end{proof}

\section{References}

For an overview of numerical eigenvalue methods, see the monograph by Saad
\cite{34}. The QR algorithm and practical implementation details, including the
Hessenberg reduction and tridiagonal and bidiagonal reduction, can be found in Demmel
\cite{13} and in Trefethen and Bau \cite{43}.

The general form of Gershgorin's theorem states that, if the union of $k$ Gershgorin
disks are disjoint from the other Gershgorin disks, then exactly $k$ of the
eigenvalues of the matrix are in the union of those $k$ Gershgorin disks. This
version of Gershgorin's theorem and extensions of it can be found in Varga \cite{44}.

The matrix
\[
B_N \;=\; \begin{pmatrix}
2 & -1 & \dots & 0\\
-1 & \ddots & \ddots & \vdots\\
\vdots & \ddots & \ddots & -1\\
0 & \dots & -1 & 2
\end{pmatrix}
\]
appears in the finite difference discretization of the one dimensional Poisson
equation. It has many interesting properties, which have a special symmetry about
them; details can be found in the section ``My Favorite Matrix'' from Gil Strang's
``Essays in Linear Algebra'' \cite{42}.

An elegant complex analysis proof of the Fundamental Theorem of Algebra (which was
used to show that an $n \times n$ matrix has exactly $n$ eigenvalues, when counted
with their multiplicities) is included, with detailed historical notes, Lax and
Zalcman \cite{27}.
